% \chapter{GIỚI THIỆU TỔNG QUAN VỀ ĐỀ TÀI}

% \section{Bối cảnh hoạt động tín dụng và rủi ro tín dụng tại Việt Nam}

% \quad Hoạt động tín dụng có vai trò đặc biệt quan trọng trong hệ thống ngân hàng thương mại Việt Nam. Nguồn thu từ tín dụng vẫn chiếm tỷ trọng lớn trong tổng thu nhập hoạt động của nhiều ngân hàng. Tuy nhiên, đi cùng với tăng trưởng tín dụng là áp lực kiểm soát rủi ro, đặc biệt trong các giai đoạn nền kinh tế chịu tác động bởi dịch bệnh, biến động lãi suất, suy giảm thị trường bất động sản, khó khăn của doanh nghiệp và sự thay đổi trong khả năng trả nợ của khách hàng cá nhân \cite{sbv2021tt11}.

% Rủi ro tín dụng phát sinh khi khách hàng không thực hiện hoặc không có khả năng thực hiện đầy đủ nghĩa vụ trả nợ theo hợp đồng đã ký kết với ngân hàng. Rủi ro này có thể đến từ nhiều nguyên nhân như năng lực tài chính yếu kém của khách hàng, biến động thu nhập, quản trị doanh nghiệp chưa hiệu quả, thông tin tín dụng không đầy đủ, biến động kinh tế vĩ mô hoặc sự thay đổi trong môi trường kinh doanh.

% Trong giai đoạn 2020 - 2025, hệ thống ngân hàng Việt Nam chịu nhiều áp lực trong kiểm soát chất lượng tín dụng. Đại dịch Covid-19 làm suy giảm hoạt động sản xuất kinh doanh của nhiều doanh nghiệp và ảnh hưởng trực tiếp đến thu nhập của cá nhân, hộ gia đình. Sau đại dịch, nền kinh tế tiếp tục đối mặt với các vấn đề như thị trường bất động sản phục hồi chậm, trái phiếu doanh nghiệp gặp khó khăn, lãi suất biến động và áp lực nợ xấu có xu hướng gia tăng.

% \begin{table}[H]
% \centering
% \caption{Tỷ lệ nợ xấu nội bảng của hệ thống ngân hàng Việt Nam giai đoạn 2020--2025}
% \label{tab:no-xau-2020-2025}
% \begin{tabular}{|c|c|}
% \hline
% \textbf{Năm} & \textbf{Tỷ lệ nợ xấu nội bảng tham khảo (\%)} \\
% \hline
% 2020 & 1,69 \\
% \hline
% 2021 & 1,90 \\
% \hline
% 2022 & 2,00 \\
% \hline
% 2023 & 2,90 \\
% \hline
% 2024 & 3,10 \\
% \hline
% 2025 & 3,00 \\
% \hline
% \end{tabular}
% \end{table}

% \begin{figure}[H]
% \centering
% % \includegraphics[width=0.8\textwidth]{figures/no_xau_2020_2025.png}
% \caption{Xu hướng tỷ lệ nợ xấu nội bảng hệ thống ngân hàng Việt Nam giai đoạn 2020--2025}
% \label{fig:no-xau-2020-2025}
% \end{figure}

% Sự gia tăng áp lực nợ xấu cho thấy nhu cầu cấp thiết trong việc nâng cao năng lực đánh giá và dự báo rủi ro tín dụng. Thay vì chỉ đánh giá khách hàng tại thời điểm cấp tín dụng, ngân hàng cần có hệ thống xếp hạng tín dụng có khả năng cập nhật, theo dõi và cảnh báo sớm sự thay đổi trong mức độ rủi ro của khách hàng.

% \section{Tổng quan về xếp hạng tín dụng}

% \quad Xếp hạng tín dụng là quá trình đánh giá mức độ tín nhiệm và khả năng thực hiện nghĩa vụ tài chính của khách hàng vay vốn. Kết quả xếp hạng thường được biểu diễn dưới dạng điểm số, hạng tín dụng hoặc nhóm rủi ro. Dựa trên kết quả này, ngân hàng có thể phân loại khách hàng theo mức độ rủi ro khác nhau, từ đó đưa ra quyết định tín dụng phù hợp.

% Trong hoạt động ngân hàng, xếp hạng tín dụng có thể áp dụng cho nhiều nhóm khách hàng khác nhau, bao gồm khách hàng cá nhân, khách hàng hộ kinh doanh, doanh nghiệp nhỏ và vừa, doanh nghiệp lớn hoặc các định chế tài chính. Đối với mỗi nhóm khách hàng, ngân hàng thường xây dựng bộ tiêu chí đánh giá riêng nhằm phản ánh đặc điểm tài chính, hành vi trả nợ và mức độ rủi ro của từng phân khúc.

% Xếp hạng tín dụng có vai trò quan trọng trong nhiều khâu của quy trình tín dụng. Trước khi cấp tín dụng, kết quả xếp hạng giúp ngân hàng đánh giá khả năng trả nợ của khách hàng và quyết định có cấp tín dụng hay không. Trong quá trình phê duyệt khoản vay, xếp hạng tín dụng hỗ trợ xác định hạn mức, lãi suất, tài sản bảo đảm và các điều kiện tín dụng. Sau khi khoản vay được giải ngân, kết quả xếp hạng có thể được sử dụng để giám sát chất lượng tín dụng, cảnh báo sớm rủi ro và phân loại nợ.

% \section{Chuyển đổi số trong ngành ngân hàng và vai trò của dữ liệu}

% \quad Chuyển đổi số đang trở thành xu hướng tất yếu trong ngành ngân hàng. Quá trình này không chỉ thể hiện ở việc số hóa kênh giao dịch với khách hàng mà còn bao gồm việc tái cấu trúc quy trình nghiệp vụ, tự động hóa vận hành, ứng dụng phân tích dữ liệu và trí tuệ nhân tạo trong quản trị.

% Trong lĩnh vực tín dụng, chuyển đổi số giúp ngân hàng thu thập và xử lý nhiều nguồn dữ liệu hơn so với phương pháp truyền thống. Ngoài dữ liệu hồ sơ vay vốn, ngân hàng còn có thể khai thác dữ liệu giao dịch, dữ liệu lịch sử thanh toán, dữ liệu từ Trung tâm Thông tin tín dụng, dữ liệu hành vi sử dụng sản phẩm dịch vụ, dữ liệu tài sản bảo đảm và các nguồn dữ liệu nội bộ khác.

% Dữ liệu trở thành yếu tố trung tâm trong quản trị rủi ro tín dụng hiện đại. Khi dữ liệu được thu thập, làm sạch, chuẩn hóa và khai thác hiệu quả, ngân hàng có thể xây dựng các mô hình dự báo có khả năng hỗ trợ ra quyết định nhanh hơn, chính xác hơn và nhất quán hơn.

% \section{Hạn chế của phương pháp xếp hạng tín dụng truyền thống}

% \quad Mặc dù hệ thống xếp hạng tín dụng truyền thống đã hỗ trợ đáng kể cho hoạt động quản trị rủi ro, nhưng trong bối cảnh dữ liệu lớn và môi trường kinh doanh biến động nhanh, các phương pháp này vẫn tồn tại một số hạn chế.

% Thứ nhất, phương pháp truyền thống thường phụ thuộc nhiều vào kinh nghiệm chuyên gia và các quy tắc nghiệp vụ cố định. Việc đánh giá khách hàng có thể chịu ảnh hưởng bởi yếu tố chủ quan, dẫn đến kết quả xếp hạng thiếu nhất quán giữa các cán bộ tín dụng hoặc giữa các đơn vị kinh doanh.

% Thứ hai, nhiều mô hình xếp hạng truyền thống dựa trên giả định tuyến tính giữa các biến đầu vào và rủi ro tín dụng. Trong thực tế, khả năng trả nợ của khách hàng chịu ảnh hưởng bởi nhiều yếu tố có quan hệ phức tạp, phi tuyến và tương tác lẫn nhau.

% Thứ ba, phương pháp truyền thống gặp hạn chế khi xử lý dữ liệu lớn và dữ liệu đa dạng. Với số lượng khách hàng, khoản vay và giao dịch ngày càng tăng, việc đánh giá thủ công hoặc dựa trên bộ quy tắc cố định có thể không còn đáp ứng yêu cầu về tốc độ và độ chính xác.

% Thứ tư, khả năng cập nhật của hệ thống truyền thống thường chậm. Khi điều kiện kinh tế thay đổi hoặc hành vi khách hàng thay đổi, mô hình xếp hạng cần được điều chỉnh kịp thời.

% Thứ năm, nhiều hệ thống xếp hạng truyền thống chưa hỗ trợ tốt chức năng cảnh báo sớm. Trong khi đó, ngân hàng không chỉ cần biết khách hàng đang thuộc nhóm rủi ro nào mà còn cần dự báo khả năng khách hàng chuyển sang nhóm rủi ro cao hơn trong tương lai.

% \section{Phát biểu bài toán nghiên cứu}

% \quad Trong đề tài này, bài toán xếp hạng tín dụng được tiếp cận dưới góc độ bài toán phân loại có giám sát. Dựa trên dữ liệu lịch sử của khách hàng tại một ngân hàng thương mại ở Việt Nam, mô hình trí tuệ nhân tạo được huấn luyện để dự đoán hạng tín dụng hoặc nhóm rủi ro tín dụng của khách hàng.

% Giả sử tập dữ liệu nghiên cứu gồm $n$ khách hàng hoặc khoản vay, được biểu diễn dưới dạng:

% \[
% X = \{x_1, x_2, ..., x_n\}
% \]

% Trong đó, mỗi quan sát $x_i$ bao gồm các đặc trưng liên quan đến khách hàng và khoản vay, chẳng hạn như thông tin nhân khẩu học, thông tin tài chính, lịch sử tín dụng, lịch sử thanh toán, thông tin khoản vay và các biến hành vi khác.

% Biến mục tiêu của bài toán là hạng tín dụng hoặc nhóm rủi ro tín dụng, được ký hiệu là:

% \[
% Y = \{y_1, y_2, ..., y_n\}
% \]

% Trong trường hợp dữ liệu được gán nhãn theo nhóm nợ, bài toán có thể được biểu diễn dưới dạng phân loại đa lớp với:

% \[
% Y \in \{1,2,3,4,5\}
% \]

% Mục tiêu của mô hình là xây dựng hàm dự báo:

% \[
% f(X) \rightarrow Y
% \]

% sao cho kết quả dự đoán gần nhất với hạng tín dụng thực tế của khách hàng.

% \section{Phương pháp nghiên cứu}

% \quad Đề tài sử dụng kết hợp phương pháp nghiên cứu định tính và định lượng.

% Phương pháp nghiên cứu định tính được sử dụng để tổng quan cơ sở lý thuyết, phân tích các khái niệm liên quan đến tín dụng, rủi ro tín dụng, xếp hạng tín dụng và ứng dụng trí tuệ nhân tạo trong ngân hàng.

% Phương pháp nghiên cứu định lượng được sử dụng trong quá trình xử lý dữ liệu, xây dựng mô hình và đánh giá hiệu quả dự báo. Quy trình nghiên cứu định lượng bao gồm các bước chính sau:

% \begin{enumerate}
%     \item Thu thập dữ liệu tín dụng.
%     \item Làm sạch dữ liệu, xử lý giá trị thiếu, xử lý dữ liệu bất thường và chuẩn hóa định dạng dữ liệu.
%     \item Phân tích khám phá dữ liệu.
%     \item Xây dựng đặc trưng phục vụ mô hình học máy.
%     \item Chia dữ liệu thành tập huấn luyện và tập kiểm tra.
%     \item Huấn luyện các mô hình trí tuệ nhân tạo.
%     \item Đánh giá mô hình bằng các chỉ số phù hợp.
%     \item So sánh kết quả và lựa chọn mô hình có hiệu quả tốt nhất.
% \end{enumerate}

% \section{Mục tiêu nghiên cứu}

% \quad Mục tiêu tổng quát của đề tài là nghiên cứu và ứng dụng các mô hình trí tuệ nhân tạo trong xếp hạng tín dụng khách hàng dựa trên dữ liệu thực tế từ một ngân hàng thương mại ở Việt Nam.

% Để đạt được mục tiêu tổng quát, đề tài tập trung thực hiện các mục tiêu cụ thể sau:

% \begin{itemize}
%     \item Hệ thống hóa cơ sở lý thuyết về tín dụng, rủi ro tín dụng, xếp hạng tín dụng và ứng dụng trí tuệ nhân tạo trong lĩnh vực ngân hàng.
%     \item Phân tích thực trạng và vai trò của xếp hạng tín dụng trong quản trị rủi ro tại ngân hàng thương mại Việt Nam.
%     \item Thu thập, tiền xử lý và phân tích bộ dữ liệu tín dụng thực tế phục vụ xây dựng mô hình.
%     \item Xây dựng và huấn luyện một số mô hình học máy phục vụ bài toán xếp hạng tín dụng.
%     \item Đánh giá, so sánh hiệu quả của các mô hình thông qua các chỉ số như Accuracy, Precision, Recall, F1-score và Confusion Matrix.
%     \item Đề xuất mô hình phù hợp và thảo luận khả năng ứng dụng trong thực tiễn hoạt động tín dụng của ngân hàng thương mại.
% \end{itemize}


% \section{Câu hỏi nghiên cứu}

% Đề tài tập trung trả lời các câu hỏi nghiên cứu sau:

% \begin{enumerate}
%     \item Các yếu tố dữ liệu nào có ảnh hưởng đáng kể đến kết quả xếp hạng tín dụng của khách hàng tại ngân hàng thương mại?
%     \item Các mô hình trí tuệ nhân tạo có thể cải thiện hiệu quả xếp hạng tín dụng so với cách tiếp cận truyền thống hay không?
%     \item Mô hình nào cho kết quả phù hợp nhất đối với bộ dữ liệu tín dụng được nghiên cứu?
%     \item Việc ứng dụng mô hình trí tuệ nhân tạo trong xếp hạng tín dụng có thể hỗ trợ ngân hàng thương mại như thế nào trong quản trị rủi ro tín dụng?
% \end{enumerate}

% \section{Ý nghĩa khoa học và ý nghĩa thực tiễn của đề tài}

% \subsection{Ý nghĩa khoa học}

% \quad Về mặt khoa học, đề tài góp phần bổ sung cơ sở nghiên cứu về ứng dụng trí tuệ nhân tạo trong lĩnh vực xếp hạng tín dụng tại ngân hàng thương mại Việt Nam. Trong bối cảnh các nghiên cứu về trí tuệ nhân tạo trong tài chính ngày càng được quan tâm, đề tài cung cấp thêm bằng chứng thực nghiệm về khả năng áp dụng các thuật toán học máy vào bài toán đánh giá rủi ro tín dụng.

% Đề tài cũng góp phần làm rõ quy trình xây dựng mô hình xếp hạng tín dụng dựa trên dữ liệu thực tế, từ khâu tiền xử lý dữ liệu, lựa chọn đặc trưng, xử lý mất cân bằng dữ liệu, huấn luyện mô hình đến đánh giá kết quả.

% Ngoài ra, nghiên cứu giúp làm rõ sự phù hợp của một số thuật toán học máy trong bối cảnh dữ liệu tín dụng tại Việt Nam. Việc so sánh hiệu quả giữa các mô hình giúp xác định ưu điểm và hạn chế của từng phương pháp, từ đó đóng góp vào cơ sở lựa chọn mô hình phù hợp cho bài toán xếp hạng tín dụng.

% \subsection{Ý nghĩa thực tiễn}

% \quad Về mặt thực tiễn, đề tài có ý nghĩa đối với ngân hàng thương mại trong việc nâng cao chất lượng đánh giá và quản trị rủi ro tín dụng. Mô hình trí tuệ nhân tạo nếu được xây dựng phù hợp có thể hỗ trợ ngân hàng đánh giá khách hàng một cách nhanh chóng, nhất quán và dựa trên dữ liệu.

% Kết quả nghiên cứu có thể hỗ trợ cán bộ tín dụng trong quá trình thẩm định hồ sơ, giúp nhận diện khách hàng có mức độ rủi ro cao, đồng thời góp phần giảm bớt sự phụ thuộc vào đánh giá chủ quan. Bên cạnh đó, mô hình có thể được sử dụng như một công cụ cảnh báo sớm, giúp ngân hàng chủ động theo dõi các khoản vay có dấu hiệu suy giảm chất lượng.

% Đối với công tác quản trị ngân hàng, ứng dụng trí tuệ nhân tạo trong xếp hạng tín dụng giúp nâng cao năng lực quản lý danh mục tín dụng, tối ưu hóa phân bổ vốn, hỗ trợ xác định chính sách lãi suất theo mức độ rủi ro và tăng hiệu quả trích lập dự phòng.

% Đối với khách hàng, hệ thống xếp hạng tín dụng dựa trên dữ liệu có thể góp phần nâng cao tính minh bạch và khách quan trong quá trình đánh giá tín dụng. Những khách hàng có lịch sử tín dụng tốt có thể được tiếp cận vốn nhanh hơn, với điều kiện tín dụng phù hợp hơn.


\chapter{GIỚI THIỆU TỔNG QUAN VỀ ĐỀ TÀI}

\section{Bối cảnh hoạt động tín dụng và rủi ro tín dụng}

Hoạt động tín dụng có vai trò đặc biệt quan trọng đối với ngân hàng thương mại. Thông qua hoạt động cấp tín dụng, ngân hàng cung ứng vốn cho cá nhân và tổ chức, đồng thời tạo ra nguồn thu từ lãi và phí. Tuy nhiên, do nghĩa vụ trả nợ được thực hiện trong tương lai, ngân hàng luôn phải đối mặt với sự không chắc chắn liên quan đến khả năng trả nợ của khách hàng.

Rủi ro tín dụng phát sinh khi khách hàng hoặc đối tác không thực hiện đầy đủ nghĩa vụ tài chính theo cam kết. Rủi ro này có thể bắt nguồn từ năng lực tài chính suy giảm, dòng tiền không ổn định, quản trị doanh nghiệp yếu, biến động thu nhập, thay đổi môi trường kinh doanh, biến động lãi suất, thông tin không đầy đủ hoặc hành vi gian lận. Hậu quả của rủi ro tín dụng có thể thể hiện qua nợ quá hạn, nợ xấu, chi phí dự phòng tăng, lợi nhuận giảm và mức độ an toàn vốn suy giảm.

Trong quản trị hiện đại, đánh giá tín dụng không nên chỉ được thực hiện một lần tại thời điểm phê duyệt. Ngân hàng cần cập nhật thông tin trong toàn bộ vòng đời khoản vay, theo dõi sự thay đổi của dư nợ, hành vi thanh toán, thời hạn, tình trạng cơ cấu lại và nhóm nợ. Điều này tạo ra nhu cầu xây dựng các mô hình có khả năng học từ dữ liệu lịch sử và hỗ trợ dự báo rủi ro theo thời gian.

\section{Tổng quan về chấm điểm và xếp hạng tín dụng}

Chấm điểm tín dụng là quá trình sử dụng dữ liệu để tính toán một điểm số, một xác suất hoặc một giá trị định lượng phản ánh mức độ rủi ro của khách hàng. Xếp hạng tín dụng là quá trình chuyển kết quả đánh giá thành các nhóm hoặc hạng rủi ro có ý nghĩa quản trị. Hai khái niệm có liên hệ chặt chẽ nhưng không hoàn toàn đồng nhất.

Trong hoạt động ngân hàng, kết quả chấm điểm và xếp hạng có thể được sử dụng để hỗ trợ quyết định cấp tín dụng, xác định hạn mức, xây dựng chính sách lãi suất, giám sát khoản vay, phân loại khách hàng, hỗ trợ trích lập dự phòng và xây dựng hệ thống cảnh báo sớm.

Đối với nghiên cứu này, đầu ra cần dự báo là nhóm nợ từ 1 đến 5. Vì vậy, thuật ngữ ``phân loại nhóm nợ'' được sử dụng để mô tả trực tiếp bài toán, trong khi ``chấm điểm tín dụng'' và ``xếp hạng tín dụng'' được sử dụng trong phần cơ sở lý thuyết và tổng quan nghiên cứu.

\section{Chuyển đổi số trong ngân hàng và vai trò của dữ liệu}

Chuyển đổi số trong ngân hàng không chỉ là số hóa kênh giao dịch mà còn bao gồm tái cấu trúc quy trình nghiệp vụ, tự động hóa vận hành, xây dựng nền tảng dữ liệu và ứng dụng các mô hình phân tích trong quản trị.

Trong tín dụng, ngân hàng có thể khai thác nhiều nhóm dữ liệu như thông tin định danh, lịch sử hợp đồng, dư nợ, hạn mức, số ngày quá hạn, thời hạn khoản vay, mục đích vay, nhóm nợ, trạng thái cơ cấu lại và dữ liệu hành vi. Khi dữ liệu được làm sạch và quản trị phù hợp, các mô hình học máy có thể hỗ trợ nhận diện quy luật mà quy tắc thủ công khó phát hiện. Tuy nhiên, dữ liệu ngân hàng cũng có những thách thức như giá trị thiếu, sai lệch nghiệp vụ, mất cân bằng lớp, thay đổi theo thời gian và nguy cơ rò rỉ thông tin.

\section{Thành tựu và hạn chế của các phương pháp đánh giá tín dụng truyền thống}

\subsection{Phương pháp chuyên gia}

Phương pháp chuyên gia dựa trên kinh nghiệm của cán bộ tín dụng và các khung đánh giá như 5C. Thành tựu của phương pháp này là hệ thống hóa các yếu tố định tính và định lượng cần xem xét, cho phép sử dụng thông tin về uy tín, năng lực quản trị, triển vọng ngành và tài sản bảo đảm.

Tuy nhiên, phương pháp này phụ thuộc vào kinh nghiệm cá nhân, khó bảo đảm tính nhất quán, khó mở rộng và khó lượng hóa mức đóng góp của từng yếu tố.

\subsection{Altman và mô hình Z-score}

Altman \cite{altman1968financial} đã đưa phân tích phân biệt đa biến vào dự báo phá sản doanh nghiệp thông qua mô hình:
\[
Z = 1.2X_1 + 1.4X_2 + 3.3X_3 + 0.6X_4 + 1.0X_5.
\]

Thành tựu nổi bật của Altman là chứng minh rằng việc kết hợp nhiều tỷ số tài chính có thể dự báo khó khăn tài chính tốt hơn việc xem xét từng chỉ số riêng lẻ. Công trình này tạo nền tảng cho dự báo phá sản và xếp hạng doanh nghiệp.

Hạn chế của mô hình là phụ thuộc vào mẫu doanh nghiệp, ngành, chuẩn kế toán và giai đoạn nghiên cứu. Mô hình sử dụng chủ yếu dữ liệu kế toán, giả định quan hệ tuyến tính và không phù hợp trực tiếp với dữ liệu khách hàng cá nhân hoặc dữ liệu hành vi.

\subsection{Ohlson và mô hình xác suất phá sản}

Ohlson \cite{ohlson1980financial} sử dụng hồi quy logistic để ước lượng xác suất phá sản. Đóng góp quan trọng là chuyển đầu ra từ một điểm phân biệt sang xác suất, tạo điều kiện thuận lợi cho lựa chọn ngưỡng và diễn giải rủi ro.

Hạn chế là quan hệ giữa log-odds và biến đầu vào được giả định tuyến tính, còn các tương tác và quan hệ phi tuyến phải được thiết kế thủ công.

\subsection{Hand và Henley}

Hand và Henley \cite{hand1997statistical} hệ thống hóa các phương pháp phân loại thống kê trong consumer credit scoring. Hai tác giả làm rõ rằng không tồn tại một thuật toán duy nhất luôn tốt nhất trên mọi bộ dữ liệu; kết quả phụ thuộc vào chất lượng dữ liệu, lựa chọn biến, cách chia mẫu và tiêu chí đánh giá.

Hạn chế chủ yếu xuất phát từ bối cảnh thời điểm nghiên cứu: boosting, explainable AI, fairness và dữ liệu quy mô lớn chưa phát triển như hiện nay; phần lớn nội dung tập trung vào phân loại nhị phân.

\subsection{Lyn C. Thomas}

Thomas \cite{thomas2000survey} có đóng góp nổi bật trong việc hệ thống hóa application scoring và behavioural scoring. Ông mở rộng credit scoring từ quyết định phê duyệt ban đầu sang quản lý khách hàng trong suốt vòng đời khoản vay.

Hạn chế của nhiều mô hình truyền thống được Thomas tổng hợp là tập trung vào dữ liệu có cấu trúc, phân loại tốt--xấu và giả định quan hệ tương đối ổn định theo thời gian.

\subsection{Jonathan N. Crook}

Crook có nhiều đóng góp trong reject inference, behavioural scoring, survival analysis và mô hình rủi ro động. Các nghiên cứu của ông làm rõ selection bias do dữ liệu huấn luyện thường chỉ gồm khách hàng đã được duyệt và mở rộng bài toán từ ``có vỡ nợ hay không'' sang ``khi nào vỡ nợ''.

Hạn chế của reject inference là không thể quan sát chắc chắn nhãn thật của hồ sơ bị từ chối; survival model cũng phụ thuộc vào các giả định có thể không phù hợp với mọi dữ liệu.

\subsection{Baesens và Lessmann}

Baesens và Lessmann có đóng góp quan trọng trong benchmark giữa các mô hình thống kê và học máy. Lessmann và cộng sự \cite{lessmann2015benchmarking} cho thấy thứ hạng mô hình thay đổi theo dữ liệu, thước đo và quy trình thực nghiệm; ensemble thường cạnh tranh nhưng không thể mặc định luôn tối ưu.

Hạn chế của nhiều benchmark là sử dụng dữ liệu công khai nhỏ, cũ, chủ yếu nhị phân và chưa phản ánh đầy đủ dữ liệu ngân hàng thực tế.

\subsection{Nhận xét tổng hợp}

Các công trình nền tảng đã chuyển đánh giá tín dụng từ chủ quan sang định lượng, tạo ra điểm số và xác suất rủi ro, hình thành application scoring, behavioural scoring và quy trình kiểm định mô hình. Tuy nhiên, các hạn chế chung còn lại gồm tập trung vào nhị phân, khó mô hình hóa phi tuyến, phụ thuộc vào xây dựng biến thủ công, chưa xử lý đầy đủ mất cân bằng, thứ tự, drift và explainability.

\section{Phát biểu bài toán nghiên cứu}

Cho tập dữ liệu:
\[
\mathcal{D}=\{(\mathbf{x}_i,y_i,t_i)\}_{i=1}^{N},
\]
trong đó $\mathbf{x}_i$ là vector đặc trưng, $y_i\in\{1,2,3,4,5\}$ là nhóm nợ và $t_i$ là thời điểm ghi nhận.

Mục tiêu là xây dựng:
\[
f_{\theta}:\mathbf{x}_i\longrightarrow(p_{i1},p_{i2},p_{i3},p_{i4},p_{i5}),
\]
với:
\[
p_{ik}=P(Y_i=k\mid\mathbf{x}_i),\qquad \sum_{k=1}^{5}p_{ik}=1.
\]
Nhãn dự báo:
\[
\hat{y}_i=\arg\max_{k\in\{1,2,3,4,5\}}p_{ik}.
\]

Do năm nhóm có quan hệ về mức độ rủi ro:
\[
1<2<3<4<5,
\]
bài toán cần vừa đạt hiệu quả tổng thể, vừa nhận diện nhóm rủi ro cao, giảm khoảng cách sai lệch, hạn chế đánh giá thấp rủi ro, duy trì khả năng tổng quát hóa theo thời gian và bảo đảm khả năng giải thích.

\subsection{Điểm phát triển so với các công trình đi trước}

Nghiên cứu phát triển hơn các công trình đi trước ở các khía cạnh:
\begin{enumerate}
    \item chuyển từ nhị phân sang phân loại trực tiếp năm nhóm nợ;
    \item xem xét thứ tự của nhãn bằng $|y_i-\hat{y}_i|$;
    \item sử dụng dữ liệu thực tế đã ẩn danh từ ngân hàng Việt Nam;
    \item xây dựng đặc trưng nghiệp vụ về thời hạn, mục đích vay, cơ cấu nợ và tính nhất quán;
    \item ưu tiên chia dữ liệu theo thời gian;
    \item xem mất cân bằng lớp là thành phần trung tâm;
    \item so sánh sáu mô hình trong cùng một quy trình;
    \item kết hợp hiệu quả dự báo với feature importance và SHAP;
    \item lựa chọn mô hình theo hiệu quả, nhận diện rủi ro, sai lệch thứ tự và khả năng giải thích.
\end{enumerate}

Tính mới của luận văn không nằm ở việc phát minh một thuật toán mới, mà nằm ở cách xây dựng, đánh giá và tổ chức bài toán phù hợp hơn với dữ liệu tín dụng thực tế.

\section{Phương pháp nghiên cứu}

Luận văn kết hợp nghiên cứu định tính để tổng quan tài liệu với nghiên cứu định lượng để xử lý dữ liệu, xây dựng mô hình và đánh giá thực nghiệm. Quy trình gồm thu thập, làm sạch, xây dựng đặc trưng, EDA, chia theo thời gian, xử lý mất cân bằng, huấn luyện, đánh giá và giải thích mô hình.

\section{Mục tiêu nghiên cứu}

\subsection{Mục tiêu tổng quát}

Xây dựng và đánh giá các mô hình học máy nhằm hỗ trợ phân loại năm nhóm nợ trên dữ liệu thực tế của một ngân hàng thương mại tại Việt Nam.

\subsection{Mục tiêu cụ thể}

\begin{enumerate}
    \item Hệ thống hóa lý thuyết và nghiên cứu liên quan.
    \item Làm sạch và xây dựng đặc trưng từ dữ liệu thực tế.
    \item Xây dựng Logistic Regression làm baseline.
    \item So sánh với Decision Tree, Random Forest, XGBoost, LightGBM và CatBoost.
    \item Đánh giá ảnh hưởng của mất cân bằng lớp.
    \item Phân tích tính thứ tự và mức độ nghiêm trọng của sai số.
    \item Kiểm tra khả năng tổng quát hóa theo thời gian.
    \item Giải thích kết quả bằng các phương pháp phù hợp.
\end{enumerate}

\section{Câu hỏi nghiên cứu}

\begin{enumerate}
    \item Những đặc trưng nào ảnh hưởng đáng kể đến phân loại nhóm nợ?
    \item Học máy có cải thiện so với Logistic Regression hay không?
    \item Mô hình nào cân bằng tốt nhất giữa hiệu quả tổng thể và nhận diện nhóm rủi ro?
    \item Xử lý mất cân bằng ảnh hưởng thế nào đến nhóm 3, 4, 5?
    \item Xem xét thứ tự có giúp giảm sai số nghiêm trọng hay không?
    \item Mô hình có ổn định trên dữ liệu thời kỳ sau hay không?
    \item Các biến quan trọng có phù hợp với nghiệp vụ tín dụng hay không?
\end{enumerate}

\section{Ý nghĩa khoa học và thực tiễn}

Về khoa học, luận văn bổ sung bằng chứng thực nghiệm về phân loại năm nhóm nợ, kết hợp mất cân bằng, yếu tố thời gian và khả năng giải thích. Về thực tiễn, kết quả có thể hỗ trợ chuẩn hóa đánh giá, nhận diện khoản vay suy giảm và xây dựng cảnh báo sớm. Mô hình chỉ đóng vai trò hỗ trợ, không thay thế hoàn toàn quyết định của cán bộ tín dụng.
